% page 433 of the facsimile (image p-432 of the scan)
% block 165.1 x 242.0 mm ; left margin 36.9 mm
\pgc{15.240mm}{0.000mm}{
\l{\cel{52}425}
\l{}
\l{\sou{pasamano}.\cel{2}Balustrado qua garnisas la bordo di eskalero, uzebla}
\l{\cel{3}kom apogo-punto da la personi qui acensas o decensas. - IS.}
\l{}
\l{\sou{pasero}.\cel{3}Ucelo\cel{2}kono-beka, kun plumaro griza, tre ordinara en}
\l{\cel{3}west-Europa. - L.\cel{1}\sou{passer}.\cel{3}FIL}
\l{}
\l{\sou{\textquotedbl{}pasha\textquotedbl{}}. (\sou{olim)}. Oficiero Turka,komisita por administrar pro-}
\l{\cel{3}vinco.}
\l{}
\l{\sou{pasifloro}. (\sou{bot.)}\cel{3}Planto de Amerika, di qua la semini saporas}
\l{\cel{3}quale olti di grenado. - DEFI.}
\l{}
\l{\sou{pasiono.}\cel{2}I. Movo violentema, impetuoza, di la psiko ad to quon}
\l{\cel{3}lu deziras. - II. Amoro violentema. - DEFIS.}
\l{}
\l{\sou{pasiv}o. - I. (\sou{komerco}).To quon onu debas. - II. (\sou{gram.})\cel{2}La voco}
\l{\cel{3}qua indikas ke la subjekto subisas la ago quan expresas la}
\l{\cel{3}verbo. -\cel{1}\sou{Pasiva}\cel{1}:Qua subisas la ago (da ulu, da ulo). - DEFIRS}
\l{}
\l{\sou{pasko}. (\sou{liturgio katol}.) Festo solena, memorige di la rezurekto}
\l{\cel{3}di Kristo. - II. (\sou{liturgio judala}). Festo solena, celebrata}
\l{\cel{3}memorige di la ekiro de Egiptia. - EFIRS.}
\l{}
\l{\sou{pasmento}. Interplekturo de fili ora, arjenta, silka - facita}
\l{\cel{3}sur kuseno, per spindeli e pingli, uzata kom bordumo o kom}
\l{\cel{3}aplikajo sur vesti, mobli, e c. - DEFIRS.}
\l{}
\l{\sou{paspelo}. Bendo ek silko, drapo, per qua onu garnisas la bordo}
\l{\cel{3}di ula parti di vesto, di uniformo. - DF.}
\l{}
\l{\sou{pasporto}.\cel{2}Dokumento emisata da Stato qua igas sua regnati darfar}
\l{\cel{3}voyajar libere, e sekurigas a li lia protekteso da la repre-}
\l{\cel{3}zentisti diplomacala di ta stato en la landi cetera. - EFIRS.}
\l{}
\l{\sou{pasto}. I. Farino dilutita e petrisita por facar pano, galeto,}
\l{\cel{3}e c. - II. Nomo di preparajo diversa en sukrajerii, apoteki,}
\l{\cel{3}parfumifeyi, e c. - DEFIS.}
\l{}
\l{\sou{pastelo.}\cel{2}Krayono ek farbi triturita e pastigita per gum-aquo.}
\l{\cel{3}DEFIRS.}
\l{}
\l{\sou{pasterno}. Infra parto di la gambo di kavalo o di bovo, ube}
\l{\cel{3}onu fixigas lige la entravilo qua retenas lu kande lu}
\l{\cel{3}pasturas. -\cel{2}EFI.}
\l{}
\l{\sou{pasteto}. Kukajo qua kontenas karno, fisho, e c., koquita en}
\l{\cel{3}pasto. - DEFIRS.}
\l{}
\l{\sou{pasteurisar}.\cel{2}(\sou{trans.)}\cel{4}Steriligar (liquido) per varmigar lu}
\l{\cel{3}til (segun lua naturo) 55o/75o C. dum kelka tempo, e koldigar}
\l{\cel{3}lu bruske. (Onu pasteurisas lakto, kom ex.) - DEFIRS.}
\l{}
\l{pastilo.\cel{2}Bonbono qua havas la formo di mikra disko, ek sukro}
\l{\cel{3}aromatizita, ek chokolado, e c. - DEFIRS.}
}
